A validation test is a decision rule with a size and a power that the literature recommending it rarely states. We measure the error rate that three corrections remove in the validation battery of an open risk-model engine, running each rule and the rule it replaces on the same seeded draws of generators with a known truth, with the standard error of every rate and, where the null law is known, the exact size. The coverage test of a split-conformal set reads the covered count against the Beta-Binomial law that a finite calibration window implies; the resilience test reads the AUC drop on the worst subsample against a permutation reference built from development rows; the situation-testing light sets its amber cut from a measured size and power. At a test window as long as the calibration window the binomial coverage test rejects an exchangeable window in \CovSinBinB{} of the replications (standard error \CovSeBinB{}) against \CovSinBBB{} for the Beta-Binomial test, whose exact size never exceeds \CovBBMaxExact{}; a binomial test with its cut calibrated to the exact size has the same power within \ExPowGapMax{}. With 500 evaluation rows a fixed drop of 0.05 raises a false alarm in \ShRateFixedFiveCZ{} of shift-free replications against \ShRatePermCZ{} for the permutation reference, which detects a shift of two standard deviations in \ShRatePermCT{} against \ShRateFixedFiveCT{}. With 400 applicants the amber cut 0.20 raises a false alarm in \SiRateCutTwentySmallZ{} against \SiRateCutTenSmallZ{} for the cut 0.10, at a cost in power at small gaps. Of \HypTotal{} predictions, \HypRefuted{} are refuted: the size of an amber cut depends on the number of applicants and on the data-generating process.
